\scanpage{234}\mspage{89}
pratique, ce sont \unsure{les|ces} stabiliseurs, plus que les
transporteurs $\mathbb L_W(A)$, qui sont importants (quand
les deux ne coïncident pas). C'est ainsi que
\inserted{dans $\operatorname{Cat}$, si $W$ est un local. fond.,}
les morphismes $W$-lisses (resp\textsuperscript{\unsure{t}} $W$-propres)
$\operatorname{Liss}_W$, $(\operatorname{Prop}_W)$, qui sont les stabiliseurs des
ensembles $W_{\mathrm{as}}$ $(W_{\mathrm{cas}})$ des \unsure{ens} de morphismes
$W$-asphériques (resp. $W$-coasphériques),
jouent un rôle essentiel alors que je ne
sais pas caractériser de façon sensible les
flèches des transporteurs $\mathbb L_W(W_{\mathrm{as}})$, $\mathbb L_W(W_{\mathrm{cas}})$,
lesquels contiennent, en plus des morphismes
$W$-lisses (resp. $W$-propres) \deleted{les} $W$-fibrations
$\operatorname{Fib}_W$. Je ne sais pas \inserted{comment} caractériser une
classe de morphismes qui apparaît aussi
hétéroclite ! De même, il y a lieu
d'introduire et de caractériser le
stabiliseur de $W_{\mathrm{bias}}$ (ens. des flèches
$W$-biasphériques), comme les morphismes
``$W$-nets'' $\operatorname{Net}_W$, \deleted{\unsure{par}} une description
similaire à celle des $\operatorname{Liss}_W$, $\operatorname{Prop}_W$,
utilisant les bilocalisés $\underset{u}{{}_{\backslash}X_{/}}$, $\underset{\unsure{v|u}}{{}_{\backslash}Y_{/}}$ au lieu
des localisés $X_{/x},Y_{/y}$ ou des colocalisés
$\,{}_{x\backslash}X,{}_{y\backslash}Y$. Il y a lieu enfin de déterminer
le stabiliseur de l'ensemble $\unsure{\operatorname{Equ}|\operatorname{Eqv}}$ des
$\operatorname{Cat}$-\unsure{équivalences} (ce qui ne dépend pas de la
donnée d'un $W$\unsure{)}, -- mais on a $\operatorname{Equ}\subset\unsure{\mathfrak{P}|\mathcal{P}}(W)$
\inserted{pour tt $W$ fondamental}).
% archive sheet number: 234
